
\setcounter{footnote}{16}

\SGARefTarget{sga3:viii:section:5}\section*{5. Quotient of an Affine Scheme by a Diagonalizable Group Acting Freely}
\addcontentsline{toc}{section}{5. Quotient of an Affine Scheme by a Diagonalizable Group Acting Freely}

\noindent\footnote{Editor's note: We have added the following sentence.}
We denote by \(\mathcal M\) the set of faithfully flat quasi-compact
morphisms, and recall that torsors are understood with respect to the
\((\mathrm{fpqc})\) topology.

\medskip
\noindent\SGARefTarget{sga3:viii:theorem:5:1}\textbf{Theorem 5.1.}
\begin{itshape}
Let \(S\) be a prescheme, \(M\) an ordinary commutative group,
\(G=D_S(M)\) the diagonalizable group over \(S\) that it defines, and
\(P\) an \(S\)-prescheme affine over \(S\), on which \(G\) acts freely on
the right.

Then the equivalence relation defined by \(G\) on \(P\) is
\(\mathcal M\)-effective (cf. IV,
\SGARefLink{sga3:iv:subsection:3:4}{3.4}); that is, the quotient
\(X=P/G\) exists and \(P\) is a torsor over \(X\) under
\(G_X=D_X(M)\). Moreover, \(P/G\) is affine over \(S\). More precisely,
if \(P\) is defined by the \(M\)-graded algebra \(\mathcal A\), then
\(P/G\) is isomorphic to \(\operatorname{Spec}(\mathcal A_0)\), where
\(\mathcal A_0=\mathcal A^G\) is the degree-zero component of
\(\mathcal A\).
\end{itshape}

\smallskip
\noindent\textit{Proof.}
Put \(X=\operatorname{Spec}(\mathcal A_0)\). The inclusion
\(\mathcal A_0\to\mathcal A\) gives a natural morphism \(P\to X\), which
is invariant under the action of \(G\). Thus \(P\) becomes an
\(X\)-prescheme, affine over \(X\), with operator group
\(G_X=D_X(M)\); the hypothesis that \(G\) acts freely on \(P/S\) implies
that \(G_X\) acts freely on \(P/X\). It remains only to show that \(P\)
is a principal homogeneous bundle under \(G_X\), using the fact that
\(\mathcal B_0=\mathcal O_X\), where \(\mathcal B\) is the \(M\)-graded
algebra over \(X\) defining \(P/X\).

We may therefore suppose that \(X=S\) and that \(S\) is affine. Then
\(P\) is affine, defined by an \(M\)-graded ring \(A\); write \(A_m\) for
its homogeneous part of degree \(m\), so that
\(S=\operatorname{Spec}(A_0)\). In view of
\SGARefLink{sga3:viii:proposition:4:6}{4.6}, it remains to verify
\[
  A_m\cdot A_{-m}=A_0
  \qquad\text{for every \(m\in M\).}
  \tag{x}\label{sga3:viii:theorem:5:1:eq:x:01}
\]
A direct calculation also shows that
\SGARefLink{sga3:viii:theorem:5:1:eq:x:01}{(x)} is equivalent to saying
that
\[
  P\times_SG\longrightarrow P\times_SP
\]
is a closed immersion, rather than merely a monomorphism (as follows
from the hypothesis that \(G\) acts freely); equivalently, the
homomorphism of affine rings

\[
  \theta:A\otimes_{A_0}A\longrightarrow A(M)
\]
is surjective.\footnote{Editor's note: Indeed, for every \(b\in A\) and
\(a_m\in A_m\), one has
\(\theta(b\otimes a_m)=ba_m\otimes e_m\); hence the surjectivity of
\(\theta\) is equivalent to \SGARefLink{sga3:viii:theorem:5:1:eq:x:01}{(x)}.}
This proves \SGARefLink{sga3:viii:theorem:5:1}{5.1} if one assumes that
the equivalence relation defined by \(G\) on \(P\) is closed. We shall
show, however, that this hypothesis already follows from the fact that
\(G\) acts freely. This is also implicit in
\SGARefLink{sga3:viii:theorem:5:1}{Theorem 5.1}: the scheme
\(G\times_SP\) must be isomorphic to \(P\times_XP\), which is closed in
\(P\times_SP\) because \(X\) is affine, hence separated, over \(S\).

Let \(R=P\times_SG\). The hypothesis that \(G\) acts freely, that is,
that \(R\to P\times_SP\) is a monomorphism, says that the diagonal
morphism
\[
  R\longrightarrow
  R\times_{(P\times_SP)}R=R'
\]
is an isomorphism. We have \(R=\operatorname{Spec}(A(M))\) and
\[
  R'=\operatorname{Spec}\bigl(A(M\times M)/K\bigr),
\]
\footnote{Editor's note: We have denoted by \(K\) the ideal denoted by
\(J\) in the original, in order to distinguish it from the ideals
\(J_p\) of \(A_0\) occurring in
\SGARefLink{sga3:viii:theorem:5:1:eq:xxx:01}{(xxx)}. We have also made
explicit below that the relations
\SGARefLink{sga3:viii:theorem:5:1:eq:xx:01}{(xx)} are equivalent to the
equality \(K=K'\); see what follows.}
where \(K\) is the ideal generated by the elements
\[
  x_m(e_{m,0}-e_{0,m}),
  \qquad m\in M,\quad x_m\in A_m.
\]
Let
\[
  \phi:A(M\times M)\longrightarrow A(M)
\]
be the surjective ring homomorphism defined by
\[
  x e_{m,n}\longmapsto x e_{m+n}
  \qquad (m,n\in M,\ x\in A),
\]
where the \(e_m\), respectively
\(e_{m,n}=e_m\otimes e_n\), are the elements of the canonical basis of
\(A(M)\), respectively \(A(M\times M)\). The diagonal morphism
\(R\to R'\) then corresponds to the homomorphism
\[
  \bar\phi:A(M\times M)/K\longrightarrow A(M)
\]
obtained by passing to the quotient by \(K\). The kernel of \(\phi\) is
the ideal \(K'\) generated by
\[
  d_m=e_{m,0}-e_{0,m}.
\]
\emph{[The source prints \(\ker(\bar\phi)=K'\); after quotienting by
\(K\), the corresponding kernel is \(K'/K\).]}
We have \(K\subseteq K'\), and the hypothesis that \(G\) acts freely on
\(P\), that is, that \(\bar\phi\) is an isomorphism, is equivalent to
\(K'=K\), expressed by the relations
\[
  d_m\in K=\sum_p A(M\times M)\,A_p\,d_p
  \qquad\text{for every \(m\in M\).}
  \tag{xx}\label{sga3:viii:theorem:5:1:eq:xx:01}
\]

Using the natural trigrading of \(A(M\times M)\), and the fact that the
first degree of \(d_m\) is zero, this says that \(d_m\) can be written as
a sum of elements of the form
\[
  f e_{r,s}(e_{p,0}-e_{0,p}),
  \qquad f\in A_{-p}\cdot A_p.
\]
Since the total degree of \(d_m\) is \(m\), it is enough to use terms
satisfying
\[
  r+s+p=m.
\]

Thus, for every \(m\in M\), there must be an expression
\[
  \left\{
  \begin{aligned}
    d_m=e_{m,0}-e_{0,m}
      &=\sum_{r,s}\lambda_{r,s}
        (e_{m-s,s}-e_{r,m-r}),\\
    \lambda_{r,s}&\in J_p=A_p\cdot A_{-p}\subseteq A_0,
      \qquad p=m-(r+s).
  \end{aligned}
  \right.
  \tag{xxx}\label{sga3:viii:theorem:5:1:eq:xxx:01}
\]
We must deduce \SGARefLink{sga3:viii:theorem:5:1:eq:x:01}{(x)}, that is,
the relations
\[
  J_n=A_0
  \qquad\text{for every \(n\in M\).}
  \tag{xxxx}\label{sga3:viii:theorem:5:1:eq:xxxx:01}
\]
For this it is enough to establish \SGARefLink{sga3:viii:theorem:5:1:eq:xxxx:01}{the same relation} modulo every
maximal ideal of \(A_0\). Since all the hypotheses are preserved under
such a reduction, we may assume that \(A_0\) is a field.

\medskip
\noindent\SGARefTarget{sga3:viii:lemma:5:2}\textbf{Lemma 5.2.}
\textit{Under the preceding conditions, with \(A_0\) a field, if
\(M\ne0\), there exists \(p\in M-\{0\}\) such that \(J_p=A_0\).}

Indeed, otherwise the sum on the right-hand side of
\SGARefLink{sga3:viii:theorem:5:1:eq:xxx:01}{(xxx)} would vanish for
every \(m\in M\), which is absurd.

\medskip
\noindent\SGARefTarget{sga3:viii:corollary:5:3}\textbf{Corollary 5.3.}
\textit{Under the preceding conditions, but without assuming that
\(A_0\) is a field, there are finitely many elements
\(p_i\in M-\{0\}\) such that}
\[
  \sum_iJ_{p_i}=A_0.
\]

Indeed, apply \SGARefLink{sga3:viii:lemma:5:2}{5.2} to the situations
obtained from \(A/A_0\) by reduction modulo the maximal ideals of
\(A_0\).\footnote{Editor's note: It follows from
\SGARefLink{sga3:viii:lemma:5:2}{5.2} that
\(\sum_{p\ne0}J_p=A_0\); hence \(1\) can be written as a finite sum
\(\sum_i x_i\), with \(x_i\in J_{p_i}\).}

\medskip
\noindent\SGARefTarget{sga3:viii:corollary:5:4}\textbf{Corollary 5.4.}
\textit{Suppose again that \(A_0\) is a field. For every subgroup
\(N\) of \(M\) with \(N\ne M\), there exists \(p\in M-N\) such that
\(J_p=A_0\).}

Indeed, let \(M'=M/N\), and consider the \(M'\)-graded ring \(A'\), whose
underlying ring is \(A\) and whose grading is given by
\[
  A'_{m'}=\coprod_{m\in h^{-1}(m')}A_m,
\]
where \(h:M\to M'=M/N\) is the canonical homomorphism. Geometrically,
this construction amounts to considering the subgroup
\(G'=D_S(M')\) of \(G\), and the structure on \(P\) of a scheme with
operator group \(G'\) induced by the action of \(G\). It is then clear
that \(G'\) acts freely on \(P\)
\emph{[the source prints \(P'\)]}; that is, the pair \((M',A')\)
satisfies the hypotheses of \SGARefLink{sga3:viii:corollary:5:3}{5.3}.
Thus
\[
  1=\sum_i f_i g_i,
  \qquad
  f_i\in A'_{m_i'},\quad
  g_i\in A'_{-m_i'},\quad
  m_i'\in M'-\{0\}.
\]
Taking the components of the right-hand side along \(A_0\), and using
the fact that \(A_0\) is a field, immediately gives
\SGARefLink{sga3:viii:corollary:5:4}{5.4}.

Now observe that
\[
  J_pJ_q\subseteq J_{p+q}
  \qquad\text{and}\qquad
  J_p=J_{-p}.
\]

Thus, if \(N\) denotes the set of \(m\in M\) such that \(J_m=A_0\)
\emph{[the source prints \(J_p=A_0\)]}, then
\(N\) is a subgroup of \(M\). By
\SGARefLink{sga3:viii:corollary:5:4}{5.4}, it equals \(M\). This
completes the proof of \SGARefLink{sga3:viii:theorem:5:1}{Theorem 5.1}.

As noted in the course of the proof,
\SGARefLink{sga3:viii:theorem:5:1}{Theorem 5.1} implies:

\medskip
\noindent\SGARefTarget{sga3:viii:corollary:5:5}\textbf{Corollary 5.5.}
\textit{Under the hypotheses of
\SGARefLink{sga3:viii:theorem:5:1}{5.1}, the graph morphism}
\[
  P\times_SG\longrightarrow P\times_SP
\]
\textit{is a closed immersion.}

It follows immediately that:

\medskip
\noindent\SGARefTarget{sga3:viii:corollary:5:6}\textbf{Corollary 5.6.}
\textit{Let \(\sigma\) be a section of \(P\) over \(S\). Then the
morphism \(g\mapsto\sigma\cdot g\) from \(G\) to \(P\) defined by
\(\sigma\) is a closed immersion.}

\medskip
\noindent\SGARefTarget{sga3:viii:corollary:5:7}\textbf{Corollary 5.7.}
\textit{Let \(G\) and \(H\) be two \(S\)-groups, with \(G\)
diagonalizable and \(H\) affine over \(S\), and let \(u:G\to H\) be a
homomorphism of \(S\)-groups that is a monomorphism. Then \(u\) is a
closed immersion, \(H/G=X\) exists, \(H\) is a principal homogeneous
bundle over \(X\) under \(G_X\), and \(X\) is affine over \(S\).}

\medskip
\noindent\SGARefTarget{sga3:viii:corollary:5:8}\textbf{Corollary 5.8.}
\textit{Under the hypotheses of
\SGARefLink{sga3:viii:theorem:5:1}{5.1}, if \(P\) is of finite type
(respectively, of finite presentation) over \(S\), then the same is
true of \(X=P/G\).}

Indeed, the hypothesis implies that the fibers of \(G_X\) are of finite
type. Hence \(G_X\) is of finite presentation over \(X\), by
\SGARefLink{sga3:viii:proposition:2:1:item:b:01}{2.1(b)}; since \(P\) is a torsor
under \(G_X\), it is therefore of finite presentation over
\(X\).\footnote{Editor's note: Cf. EGA II, 2.7.1(vi).}
It is also faithfully flat over \(X\), so the conclusion follows from
Expos\'e V, Proposition \SGARefLink{sga3:v:proposition:9:1}{9.1}.
